\scanpage{11}
De même, on aura (prop. 2, cor)
\[
\varphi_{\mathcal L}(xy)=(\varphi_{\mathcal L}(x)*\varphi_{\mathcal L}(y))\deleted{\unsure{\mathcal L}}\cdot(-1)^n\omega_A
\]
d'où en appliquant $\operatorname{ch}$ et $\unsure{\text{a)}|\alpha)}$
\[
\psi_D(\operatorname{ch}(xy))=\deleted{\unsure{\psi_D(}}\psi_D(\operatorname{ch}x)\psi_D(\operatorname{ch}y))T(-1)^n\deleted{\unsure{\exp}}\operatorname{ch}\omega_A
\]
donc comme $\operatorname{ch}(xy)=\operatorname{ch}x\operatorname{ch}y$
\[
\boxed{\psi_D(ab)=\psi_D(a)\psi_D(b)(-1)^nT\operatorname{ch}\omega_A}
\qquad\begin{cases}T=\operatorname{Todd}(t_A)\\\omega_A=\Lambda^n t_A^\vee\end{cases}
\]
\margin{Ces formules sont valables dès que $a,b\in\unsure{\operatorname{Gr}^\bullet(A)_{\mathbf Q}|\operatorname{Gr}(A)_{\mathbf Q}}$.}
Ce sont les formules \deleted{\unsure{différentielles}}\inserted{démontrées} (\unsure{cas spéciaux})\\
par Kleiman dans ``10 Exposés ...'', sur un\\
corps de base et en travaillant avec une\\
théorie cohomologique plutôt qu'avec $\operatorname{Gr}_{\mathbf Q}$, ---\\
à cela près que Kleiman\deleted{\unsure{n}} n'a pas tenu\\
compte du facteur $(-1)^n$ !

4) Il résulte du Th. de dualité pour\\
$A\times_S B\longrightarrow B$ qu'il y a un isom. fonctoriel
\[
\Phi_{\mathcal L}(\check F\otimes\omega_A[\unsure{n|-n}])\simeq\Phi_{\mathcal L}(F)^\vee
\]
i.e.
\[
\boxed{\Phi_{\mathcal L}(F)^\vee\simeq\Phi_{\mathcal L}(\check F.\omega_A[n]}
\]
Du point de vue virtuel, ceci s'écrit
\[
\varphi_{\mathcal L}(x)^\vee=(-1)^n\omega_A\varphi_{\mathcal L}(\check x)
\]
5) Bien sûr, si $u$ est un autom. de $A$ et\\
$u'$ l'autom. contragrédient de $B$, on a par\\
transport de structure un isom canonique
\[
\Phi_{\mathcal L}(u.F)\simeq u'.\Phi_{\mathcal L}(\unsure{f|F})
\]
En particulier, $\Phi_{\mathcal L}$ commute à la symétrie $s_A$. \unsure{Id.}\\
\inserted{\unsure{Id.}} pour $\varphi_{\mathcal L}$.
\mspage{10}
